\section{Evaluation}
\label{sec:analysis-evaluation}

The main engineering lesson is that partitioning, deployment, and hardening
cannot be optimised independently. The split point determines
activation-transfer size and compute placement. Kubernetes and AKS add
platform-specific communication and scheduling effects. Security hardening
then compounds the cost of the selected service topology. This means that a
decomposition that is acceptable without mTLS or confidential execution may
become less attractive once every service boundary is protected.

For the evaluated workload, \texttt{chain\_2svc} is the most defensible
non-monolithic deployment point. It avoids the excessive boundary count of
deeper chains, preserves a meaningful architectural split, and remains the
lowest-overhead chained condition in both local Kubernetes and AKS. It also
provides the practical base for RQ2: the mTLS/AuthZ hardening bundle is
measurable but bounded on this topology, and selective TEE can protect the
downstream service at a lower cost than full-chain confidential execution.

Across the evaluated stack, the monolithic Kubernetes baseline remains the
fastest path; chained deployments pay an overhead that is monotonic with
service count but bounded and non-linear; the local-to-AKS transfer is
directional rather than numerical; mTLS adds a workload-amortised but
measurable cost that grows with protected chain depth; and selective
confidential-VM execution on the downstream service adds a further bounded
cost on top of the communication-secured chain. Extending the
confidential-execution boundary to both services is treated as future work.
